\startSection
	\selectlanguage{german}
	\sectionTitle{Die Deutsche Nationalhymne}
	
	\begin{question}{Die Melodie der deutschen Nationalhymne stammt von einem bekannten Komponisten der Klassik. Wer komponierte ursprünglich diese Melodie, die heute mit dem Deutschlandlied verbunden ist?}
		\choice{Ludwig van Beethoven}
		\choice{Johann Sebastian Bach}
		\choice[correct]{Joseph Haydn}
		\choice{Richard Wagner}
		\choice{Johannes Brahms}
	\end{question}

	
	\begin{question}{Der Text des Deutschlandliedes wurde im 19. Jahrhundert verfasst und spiegelt die damaligen nationalen Ideale wider. Wer war der Autor dieses Liedtextes?}
		\choice{Johann Wolfgang von Goethe}
		\choice{Friedrich Schiller}
		\choice{Heinrich Heine}
		\choice[correct]{August Heinrich Hoffmann von Fallersleben}
		\choice{Theodor Körner}
	\end{question}

	
	\begin{question}{Das Deutschlandlied besteht aus drei Strophen, doch seit 1952 wird nur eine davon offiziell als Nationalhymne verwendet. Welche Strophe wird seitdem bei offiziellen Anlässen gesungen?}
		\choice{Die erste Strophe}
		\choice{Die zweite Strophe}
		\choice[correct]{Die dritte Strophe}
		\choice{Alle drei Strophen}
		\choice{Die erste und dritte Strophe}
	\end{question}

	
	\begin{question}{Die heute gesungene Strophe der deutschen Nationalhymne beginnt mit einem markanten Vers, der zentrale Werte der Bundesrepublik Deutschland betont. Mit welchen Worten beginnt diese Strophe?}
		\choice{\enquote{Deutschland, Deutschland über alles}}
		\choice{\enquote{Deutsche Frauen, deutsche Treue}}
		\choice[correct]{\enquote{Einigkeit und Recht und Freiheit}}
		\choice{\enquote{Blüh im Glanze dieses Glückes}}
		\choice{\enquote{Von der Maas bis an die Memel}}
	\end{question}

	
	\begin{question}{Das Deutschlandlied wurde in einer Zeit des politischen Umbruchs und der nationalen Bewegung verfasst. In welchem Jahr wurde der Text dieses Liedes ursprünglich geschrieben?}
		\choice{1797}
		\choice{1815}
		\choice[correct]{1841}
		\choice{1871}
		\choice{1918}
	\end{question}
\renderSection
